%%% Lecture 1 Slides %%%

\documentclass{beamer}

%\documentclass[handout]{beamer}
%\usepackage{pgfpages}
%\pgfpagesuselayout{2 on 1}[a4paper,border shrink=5mm]
%\pgfpagesuselayout{1 on 1}[a4paper,border shrink=5mm]



\usepackage{amssymb}
\usepackage{amsmath}
\usepackage{mathtools}
\usepackage{tikz}
\usetikzlibrary{shapes,arrows}
\usetikzlibrary{positioning}
\usetikzlibrary{trees}
\usepackage{pgfplots}
\usepackage{pdflscape}
\usepackage{graphicx}
\usepackage{colortbl}
\usepackage{listings}
\usepackage{adjustbox}

\DeclarePairedDelimiter{\ceil}{\lceil}{\rceil}

\newcommand{\bs}{\boldsymbol}
\newcommand{\mb}{\mathbf}
\newcommand{\indep}{{\bot\negthickspace\negthickspace\bot}}
\newcommand{\notindep}{{\bot\negthickspace\negthickspace\bot}{\negthinspace
  \negmedspace \negthickspace \negthickspace \diagup
  \;}}
\renewcommand{\S}{\mb{S}}
\usepackage{color}
\definecolor{direct}{HTML}{FF0000}
\definecolor{indirect}{HTML}{FF9999}
\definecolor{dirin}{HTML}{990000}
\definecolor{control}{HTML}{999999}
\definecolor{directE}{HTML}{ED1C24}
\definecolor{indirectE}{HTML}{F69679}
\definecolor{controlE}{HTML}{999999}

%\defV{{\rm Var}\,}
\def\E{{\rm E}\,}
\def\arg{{\rm arg}\,}
\def\Cov{{\rm Cov}\,}
\def\N{{\rm N}\,}
\newcommand{\plim}{{\rm plim} \,}
\DeclareMathOperator*{\argmin}{arg\,min}
\newcommand{\X}{\mathbf{X}}
\newcommand{\Y}{\mathbf{Y}}
\newcommand{\Z}{\mathbf{Z}}
\newcommand{\I}{\mathbf{I}}
\newcommand{\V}{\mathbf{\Omega}}
\newcommand{\W}{\mathbf{W}}
\newcommand{\R}{\mathbf{R}}
\newcommand{\A}{\mathbf{A}}
\newcommand{\B}{\mathbf{B}}
\newcommand{\D}{\mathbf{D}}
\newcommand{\T}{\mathbf{T}}
\newcommand{\F}{\mathbf{F}}
\renewcommand{\O}{\mathbf{O}}
\renewcommand{\H}{\mathbf{H}}
\newcommand{\g}{\mathbf{g}}
\renewcommand{\r}{\mathbf{r}}
\newcommand{\s}{\mathbf{s}}
\newcommand{\uu}{\mathbf{u}}
\newcommand{\w}{\mathbf{w}}
\newcommand{\x}{\mathbf{x}}
\newcommand{\vv}{\mathbf{v}}
\newcommand{\z}{\mathbf{z}}
\newcommand{\y}{\mathbf{y}}

\newcommand{\bi}{\begin{itemize}}
\newcommand{\ei}{\end{itemize}}
\newcommand{\be}{\begin{enumerate}}
\newcommand{\ee}{\end{enumerate}}

\newcommand{\Beta}{\boldsymbol{\beta}}
\newcommand{\btheta}{\boldsymbol{\theta}}
\newcommand{\bgamma}{\boldsymbol{\gamma}}
\newcommand{\bpi}{\boldsymbol{\pi}}
\newcommand{\arrowp}{\stackrel{p}{\rightarrow}}
\newcommand{\arrowd}{\stackrel{d}{\rightarrow}}
\newcommand{\0}{\mathbf{0}}
\newcommand{\bP}{\mathbf{P}}
\newcommand{\nn}{\nonumber}
\def\S{{\sigma}}
\def\Supp{{\rm Supp}\,}


\newcommand\independent{\protect\mathpalette{\protect\independenT}{\perp}}
\def\independenT#1#2{\mathrel{\rlap{$#1#2$}\mkern2mu{#1#2}}}

\newcommand{\vsf}{\vspace{.5cm}}

\usepackage[english]{babel}
% or whatever

\usepackage[latin1]{inputenc}
% or whatever

\usepackage{times}
\usepackage[T1]{fontenc}
\setbeamertemplate{footline}[page number]{}
\setbeamertemplate{navigation symbols}{}


\title[Lecture 14] % (optional, use only with long paper titles)
{Lecture 14}


\begin{document}

\begin{frame}
  \titlepage
\end{frame}



\frame{
\frametitle{Linear Regression Slope Coefficients}

Suppose we have de-meaned all variables.\\

\bigskip


One Covariate ($Y \sim X$):\\
$\frac{\sum\limits_{i=1}^{n} x_i y_{i}}{\sum\limits_{i=1}^{n}x_i^2}$

\bigskip
\pause

Two Covariates ($Y \sim X + Z$):\\
$\frac{\sum\limits_{i=1}^{n} \left[x_i - \left(\frac{\sum\limits_{i=1}^{n} z_i x_i}{\sum\limits_{i=1}^{n} z_i^2}\right)z_i\right] y_{i}}{\sum\limits_{i=1}^{n} \left[x_i - \left(\frac{\sum\limits_{i=1}^{n} z_i x_i}{\sum\limits_{i=1}^{n} z_i^2}\right)z_i\right]^2}$

\bigskip

\pause

A more general approach would be helpful.


}




\frame{

Define
$$
\X = \left( \begin{array}{ccccc} 1 & X_{11} & X_{21} & ... & X_{K1} \\ 1 & X_{12} & X_{22} & ... & X_{K2} \\ \vdots &  \vdots &  \vdots & \ddots &  \vdots \\  1 & X_{1n} & X_{2n} & ... & X_{Kn} \end{array} \right), \hspace{1em}  \Y = \left( \begin{array}{c} Y_1 \\ Y_2 \\ \vdots \\ Y_n \end{array} \right), \hspace{1em}  {\bf e} = \left( \begin{array}{c} e_1 \\  e_2 \\ \vdots \\  e_n \end{array} \right),% {\bf\epsilon} = \left( \begin{array}{c} \epsilon_1 \\  \epsilon_2 \\ \vdots \\  \epsilon_n \end{array} \right),
$$

$$
\text{and} \hspace{1em} \widehat \Beta =\left( \begin{array}{c} \hat \beta_0 \\ \hat \beta_1 \\  \hat \beta_2 \\ \vdots \\ \hat \beta_K \end{array} \right).
 $$
}





\frame{

Then it is equivalent to write the following:
\begin{align*}
Y_1 & = \hat \beta_0 + \hat \beta_1 X_{11} + \hat \beta_2 X_{21} + ... \hat \beta_K X_{K1} + e_1 \\
Y_2 & = \hat \beta_0 + \hat \beta_1 X_{12} + \hat \beta_2 X_{22} + ... \hat \beta_K X_{K2} + e_2 \\
& \qquad \qquad \qquad \quad  ... \\
Y_n & = \hat \beta_0 + \hat \beta_1 X_{1n} + \hat \beta_2 X_{2n} + ... \hat \beta_K X_{Kn} + e_n \\
\end{align*}

}

\frame{

Then it is equivalent to write the following:
\begin{align*}
\left( \begin{array}{c} Y_1 \\ Y_2 \\ \vdots \\ Y_n \end{array} \right) &=  \left( \begin{array}{ccccc} 1 & X_{11} & X_{21} & ... & X_{K1} \\ 1 & X_{12} & X_{22} & ... & X_{K2} \\ \vdots &  \vdots &  \vdots & \ddots &  \vdots \\  1 & X_{1n} & X_{2n} & ... & X_{Kn} \end{array} \right) \left( \begin{array}{c} \hat \beta_0 \\ \hat \beta_1 \\ \hat \beta_2 \\ \vdots \\ \hat \beta_K \end{array} \right) + \left( \begin{array}{c} e_1 \\  e_2 \\ \vdots \\  e_n \end{array} \right).
\end{align*}

}

\frame{

Then it is equivalent to write the following:
\begin{align*}
\Y & = \X \widehat \Beta + {\bf e}.
\end{align*}

}


\frame{
\frametitle{OLS Solution}

$$
\widehat \Beta  = (\X'\X)^{-1} \X' \Y
$$

}

\frame{

Residuals can be written as the following:
\begin{align*}
\left( \begin{array}{c} e_1 \\  e_2 \\ \vdots \\  e_n \end{array} \right) &= \left( \begin{array}{c} Y_1 \\ Y_2 \\ \vdots \\ Y_n \end{array} \right) -  \left( \begin{array}{ccccc} 1 & X_{11} & X_{21} & ... & X_{K1} \\ 1 & X_{12} & X_{22} & ... & X_{K2} \\ \vdots &  \vdots &  \vdots & \ddots &  \vdots \\  1 & X_{1n} & X_{2n} & ... & X_{Kn} \end{array} \right) \left( \begin{array}{c} \hat \beta_0 \\ \hat \beta_1 \\ \hat \beta_2 \\ \vdots \\ \hat \beta_K \end{array} \right).
\end{align*}
\begin{align*}
 {\bf e}  &=  \Y - \X \widehat \Beta.
\end{align*}


}


\frame{

Minimizing the sum of squared residuals:
\begin{align*}
\widehat \Beta = argmin_{\bf b} (\Y - \X \bf{b})' (\Y - \X \bf{b})
\end{align*}

Setting derivatives equal to zero and substituting $\widehat\Beta$ for $\bf b$.
\begin{align*}
2 \X^T (\Y - \X\widehat\Beta)&= 0 \\
\X^T {\bf e} &= 0
\end{align*}

Let's solve the first equation for $\widehat\Beta$.
 
}

\frame{

Recall that $\Y$ can also be written in terms of $\Beta$ and ${\bf \epsilon}$
\begin{align*}
\Y & = \X \Beta + {\bf \epsilon}.
\end{align*}

Let's draw a picture to remind ourselves of the difference between $\widehat\Beta$ and $\Beta$ and  ${\bf e}$ and ${\bf \epsilon}$.

}

\frame{
\frametitle{How to make CIs}

Two main ideas:
\begin{enumerate}
\item<1-> Normal Approximation
\begin{enumerate}
\item<2-> prove asymptotic normality of estimator
\item<3-> derive mean and variance (and SE) of estimator (why only these needed?)
\item<4-> estimate components of variance (and bias if necessary)
\item<5-> do the CI algebra
\end{enumerate}
\item<6-> Bootstrap
\end{enumerate}

}

\frame{
Bootstrapping ($\hat \beta_k$ for $k=1,...,K$)
\begin{enumerate}
\item Take a sample of size n with replacement from your sample
\item Calculate $\hat \beta_k^{(b=1)}$ from that bootstrap sample
\item Repeat for $b=1,...B$ where $B$ is big
\item  Do one of the following:
\begin{itemize}
\item Calculate bootstrapped variance $\frac{1}{B-1} \sum_{b=1}^B (\hat \beta_k^{(b)} - \overline{\hat \beta_k^{(b)}})^2$
\item Put $\{\hat \beta_k^{(b)}: b=1,...,B\}$ in order and report the $\alpha/2$ and $1-\alpha/2$ percentiles
\item ...
\end{itemize}
\end{enumerate}

%\vsf
%
%Takes care of heteroskedasticity automatically.
}

\frame{


Normal Approximation
\begin{enumerate}
\item<1-> prove asymptotic normality of estimator
\item<2-> derive mean and variance (and SE) of estimator (why only these needed?)
\item<3-> estimate components of variance (and bias if necessary)
\item<4-> do the CI algebra
\end{enumerate}

}


\frame{
\frametitle{Mean of OLS Estimator}

\begin{align*}
\widehat \Beta  &= (\X'\X)^{-1} \X' \Y\\
&= (\X'\X)^{-1} \X' (\X \Beta + {\bf \epsilon}) \\
&= (\X'\X)^{-1} \X'\X \Beta + (\X'\X)^{-1} \X'{\bf \epsilon}) \\
&= \Beta + (\X'\X)^{-1} \X'{\bf \epsilon} \\
\end{align*}

\begin{align*}
E[\widehat \Beta| \X]  &= E[\Beta| \X] + E[(\X'\X)^{-1} \X'{\bf \epsilon}| \X] \\
&= \Beta + (\X'\X)^{-1} \X'E[{\bf \epsilon}| \X] \\
&= \Beta %\because Theorem 2.2.22
\end{align*}

}

\frame{

By iterated expectations,

\begin{align*}
E[\widehat \Beta]  &=E_{\X}\{E[\widehat \Beta| \X]  \} \\
&=  E_{\X}\{ \Beta \} \\
&= \Beta
\end{align*}


}





%\frame{
%\frametitle{Consistency}
%$$
%\widehat \Beta \rightarrow_p \Beta
%$$
%
%\bigskip
%
%Handwaving Proof:
%\begin{itemize}
%\item Any of the $K$ slopes in $\Beta$ can be rewritten as a simple regression coefficient using the FWL.
%\item We've previously shown how to prove the consistency of simple regression coefficients using the LLN and CMT.
%\item There are a few details I'm leaving out here. 
%\end{itemize}
%
%}

\frame{
Sampling distributions of the simple and multiple sample slopes (in the $X$ direction).
\begin{table}[h]
\begin{center}
{\scriptsize
\begin{tabular}{r|r|r|r|r|r|rr} %cc
$i$ & \multicolumn{1}{r|}{1} & \multicolumn{1}{r|}{2}   & $\dots$ & \multicolumn{1}{r|}{12} & & & \\ % & \multicolumn{1}{c}{$MoE$}& \multicolumn{1}{c}{95\% CI}  \\
$s$  & $Y_1$, $X_1$, $Z_1$ & $Y_2$, $X_2$, $Z_2$   & $\dots$  & $Y_{12}$,$X_{12}$, $Z_{12}$  & $\hat{\beta}_1^s$ & $\hat{\beta}_1^m$\\ % & $2\cdot\sqrt{\frac{\overline{Y}_{625}(1-\overline{Y}_{625})}{625}}$ & $\overline{Y}_{625} \pm MoE$ \\
\hline
1 & 3.3,97,100 & 57.2,36,72.6   & $\dots$  & 39.3,30,40.6 & -.498 & -.499   \\ %& \visible<2->{.037} & \visible<3->{[.643,.717]}\\
2 & 30.4,54,99.3& 54.0,24,28.5  & $\dots$  & 6.2,99,100 & -.456 & -.398 \\ %& \visible<2->{.037}  & \visible<3->{[.663,.737]}\\
3 & 14.7,87,94.8 & 3.8,98,100 & $\dots$  & 3.8,98,100 & -.209 & -.190 \\ %& \visible<2->{.035} & \visible<3->{[.715,.785]}\\
4 & $\vdots$ & $\vdots$ & $\dots$  & $\vdots$ & $\vdots$ \\ %& \visible<2->{.036} & \visible<3->{[.694,.766]}\\
5 & $\vdots$ & $\vdots$  & $\dots$  & $\vdots$ & $\vdots$ \\ %& \visible<2->{.037} & \visible<3->{[.653,.727]}\\
$\vdots$ & $\vdots$  & $\vdots$ & $\vdots$ & $\vdots$  & $\vdots$ \\ % & \visible<2->{$\vdots$} & \visible<3->{$\vdots$}  \\
\hline
%\visible<1->{$f_{\cdot}(\cdot)$} & \multicolumn{1}{r|}{\visible<1->{?}} & \multicolumn{1}{r|}{\visible<1->{?}}  & $\dots$& \multicolumn{1}{r|}{\visible<1->{?}} & \visible<1->{?} \\
\visible<1->{$\E_{\cdot}(\cdot)$} & \multicolumn{1}{r|}{\visible<1->{17.3, 70.7,89.3}} & \multicolumn{1}{r|}{\visible<1->{17.3, 70.7,89.3}}  & $\dots$& \multicolumn{1}{r|}{\visible<1->{17.3, 70.7,89.3}} & \visible<1->{\textcolor{black}{-.52}} & -.33  \\
%\visible<1->{$V_{\cdot}(\cdot)$} & \multicolumn{1}{r|}{\visible<1->{436}} & \multicolumn{1}{r|}{\visible<1->{436}}  & $\dots$& \multicolumn{1}{r|}{\visible<1->{436}} & \visible<1->{\textcolor{black}{$\frac{436}{12}$}} \\
\end{tabular}}
%\caption{\emph{Sampling Distributions for $\overline{Y}_{625}$, the Margin-of-Error, and the 95\% Confidence Interval}}\label{t:samdistCI}
\end{center}
\end{table}

}



%\frame{
%\frametitle{Multivariate OVB}
%
%\begin{itemize}
%
%\item Suppose the model you've estimated is
%\begin{align*}
%\Y= \X\Beta + \epsilon
%\end{align*}
%when the model you wanted was
%\begin{align*}
%\Y= \X\Beta + \Z\Delta + \epsilon
%\end{align*}
%
%
%\end{itemize}
%}
%
%\frame{
%
%Define
%$$
%\X = \left( \begin{array}{ccccc} 1 & X_{11} & X_{21} & ... & X_{K1} \\ 1 & X_{12} & X_{22} & ... & X_{K2} \\ \vdots &  \vdots &  \vdots & \ddots &  \vdots \\  1 & X_{1n} & X_{2n} & ... & X_{Kn} \end{array} \right), \hspace{1em}   \Beta =\left( \begin{array}{c} \beta_0 \\ \beta_1 \\ \beta_2 \\ \vdots \\ \beta_K \end{array} \right),
%$$
%
%$$
%\Z = \left( \begin{array}{cccc}  Z_{11} & Z_{21} & ... & Z_{J1} \\  Z_{12} & Z_{22} & ... & Z_{J2} \\  \vdots &  \vdots & \ddots &  \vdots \\  Z_{1n} & Z_{2n} & ... & Z_{Jn} \end{array} \right), \hspace{1em}  
% \Delta =\left( \begin{array}{c}  \delta_1 \\  \delta_2 \\ \vdots \\  \delta_J \end{array} \right).
%$$
%
%%$$
%%\text{and} \hspace{1em} \Beta =\left( \begin{array}{c} \beta_0 \\ \beta_1 \\ \beta_2 \\ \vdots \\ \beta_K \end{array} \right).
%% $$
% 
% $$
%  \Y = \left( \begin{array}{c} Y_1 \\ Y_2 \\ \vdots \\ Y_n \end{array} \right), \hspace{1em}  {\bf\epsilon} = \left( \begin{array}{c} \epsilon_1 \\  \epsilon_2 \\ \vdots \\  \epsilon_n \end{array} \right).
% $$
% 
%}
%
%
%\begin{frame}
%
%\begin{itemize}
%\item Then your OLS estimator, $\hat{\Beta}$ is
%
%{\small
%\begin{align*}
%\left[\X^T \X \right]^{-1} \X^T \Y & = \hat \Beta  \\
%\left[\X^T \X \right]^{-1} \X^T [\X \Beta + \Z\Delta + {\bf\epsilon}] & = \hat \Beta \\
%\left[\X^T \X \right]^{-1} \X^T \X \Beta + \left[\X^T \X \right]^{-1} \X^T \Z \Delta + \left[\X^T \X \right]^{-1} \X^T {\bf\epsilon} & = \hat \Beta \\
%\Beta + \left[\X^T \X \right]^{-1} \X^T \Z \Delta + \left[\X^T \X \right]^{-1} \X^T {\bf\epsilon} & = \hat \Beta
%\end{align*}
%}
%
%\item Use the rules of conditional expectation to derive the bias.
%\end{itemize}
%\end{frame}
%
%\frame{
%
%
%
%
%{\small
%\begin{align*}
%\E[\Beta + \left[\X^T \X \right]^{-1} \X^T \Z \Delta + \left[\X^T \X \right]^{-1} \X^T {\bf\epsilon}] & = \E[ \hat \Beta | \X,\Z] \\
%\Beta + \left[\X^T \X \right]^{-1} \X^T \Z \Delta + \left[\X^T \X \right]^{-1} \X^T \E[{\bf\epsilon}| \X,\Z] & = \E[ \hat \Beta | \X,\Z] \\
%\Beta + \left[\X^T \X \right]^{-1} \X^T \Z \Delta  & = \E[ \hat \Beta | \X,\Z]
%\end{align*}
%}
%}
%
%
%\frame{
%\frametitle{Example with $K=2$ and $J=1$ }
%
%
%{\small
%\begin{align*}
%Z_1 &=\X \Gamma  + \nu \\
%&= \left( \begin{array}{ccc} 1 & X_{11} & X_{21}  \\ 1 & X_{12} & X_{22}  \\ \vdots   & \ddots &  \vdots \\  1 & X_{1n} & X_{2n}  \end{array} \right)  \left( \begin{array}{c} \hat \gamma_0 \\ \hat \gamma_1 \\ \hat \gamma_2 \end{array} \right)   + \left( \begin{array}{c} \nu_1 \\  \nu_2 \\ \vdots \\  \nu_n \end{array} \right) %\\
%%&= \left( \begin{array}{c} \hat \gamma_0 \\ \hat \gamma_1 \\ \hat \gamma_2 \end{array} \right) 
% \end{align*}
%}
%
% {\small
%\begin{align*}
% \left[\X^T \X \right]^{-1} \X^T \Z \Delta  & =   \left[\X^T \X \right]^{-1} \X^T Z_1  \cdot   \delta_1 \\
% & =   \hat \Gamma  \cdot   \delta_1 \\
% & =   \left( \begin{array}{c} \hat \gamma_0 \\ \hat \gamma_1 \\ \hat \gamma_2 \end{array} \right) \cdot   \delta_1 \\
% \end{align*}
%}
%
%\center
%
%Bias for $\beta_1$ is $\hat \gamma_1 \cdot   \delta_1$
%
%}
%
%\frame{
%\frametitle{Example with $K=2$ and $J=2$ (Part 1)}
%
%
%{\small
%\begin{align*}
%Z_1 &=\X \Gamma^1  + \nu^1 \\
%&= \left( \begin{array}{ccc} 1 & X_{11} & X_{21}  \\ 1 & X_{12} & X_{22}  \\ \vdots   & \ddots &  \vdots \\  1 & X_{1n} & X_{2n}  \end{array} \right)  \left( \begin{array}{c} \hat \gamma_0^1 \\ \hat \gamma_1^1 \\ \hat \gamma_2^1 \end{array} \right)   + \left( \begin{array}{c} \nu_1^1 \\  \nu_2^1 \\ \vdots \\  \nu_n^1 \end{array} \right) %\\
%%&= \left( \begin{array}{c} \hat \gamma_0 \\ \hat \gamma_1 \\ \hat \gamma_2 \end{array} \right) 
% \end{align*}
%}
%
%{\small
%\begin{align*}
%Z_2 &=\X \Gamma^2  + \nu^2 \\
%&= \left( \begin{array}{ccc} 1 & X_{11} & X_{21}  \\ 1 & X_{12} & X_{22}  \\ \vdots   & \ddots &  \vdots \\  1 & X_{1n} & X_{2n}  \end{array} \right)  \left( \begin{array}{c} \hat \gamma_0^2 \\ \hat \gamma_1^2 \\ \hat \gamma_2^2 \end{array} \right)   + \left( \begin{array}{c} \nu_1^2 \\  \nu_2^2 \\ \vdots \\  \nu_n^2 \end{array} \right) %\\
%%&= \left( \begin{array}{c} \hat \gamma_0 \\ \hat \gamma_1 \\ \hat \gamma_2 \end{array} \right) 
% \end{align*}
%}
%
%}
%
%\frame{
%\frametitle{Example with $K=2$ and $J=2$ (Part 2)}
%
%
%
% {\small
%\begin{align*}
% \left[\X^T \X \right]^{-1} \X^T \Z \Delta  & =   \left[\X^T \X \right]^{-1} \X^T Z_1  \cdot   \delta_1 +\left[\X^T \X \right]^{-1} \X^T Z_2  \cdot   \delta_2 \\
%& =   \left( \begin{array}{c} \hat \gamma_0^1 \\ \hat \gamma_1^1 \\ \hat \gamma_2^1 \end{array} \right) \cdot   \delta_1 +  \left( \begin{array}{c} \hat \gamma_0^2 \\ \hat \gamma_1^2 \\ \hat \gamma_2^2 \end{array} \right) \cdot   \delta_2 \\
% \end{align*}
%}
%
%\center
%
%Bias for $\beta_1$ is $\hat \gamma_1^1 \cdot   \delta_1 + \hat \gamma_1^2 \cdot   \delta_2$
%
%}
%
%






%\frame{
%\frametitle{Why do we care about SEs?}
%
%
%%\begin{itemize}
%%\item If x has been randomized, what do we know about the slope in the x direction for OLS with x and OLS with x and z?
%%\item<2-> Write the estimated classical SE formulas for OLS with x and OLS with x and z, where x has been randomized.
%%%\item Sampling distribution table for OLS with x and OLS with x and z, where x has been randomized
%%%\item Include formulas for SEs (or variances) and discuss bias and efficiency (including VIF and changing errors) and draw sampling distributions
%%%\item Discuss estimation of $\sigma^2$ for different models (similarities and differences).
%%%\item Present SE for heteroskedasticity with single x, and prove equals classic SE when homoskedasticity holds
%%%\item Discuss data transformation version of WLS and include in table and discussion and draw sampling distributions.
%%\end{itemize}
%
%\begin{itemize}
%\item<2-> Inference (e.g., confidence intervals)
%\item<3-> Efficiency (i.e., lower MSE)
%\end{itemize}
%
%
%
%%\visible<4>{Compare sampling distributions for OLS with one and two vars, WLS?, and Estimators of SEs on the board}
%
%}
%

\frame{
\frametitle{Law of Total Variance}


\begin{align*}
V\left[\hat \Beta\right] &= E\left[ V\left[ \hat \Beta | \X \right] \right] + V\left[ E\left[  \hat \beta  | \X   \right] \right] \\
&= E\left[ V\left[ \hat \Beta | \X \right] \right] + 0
\end{align*}

}


\frame{

We can decompose the $\hat\Beta$ into the expected value and the sampling error.

\begin{align*}
\left[\X^T \X \right]^{-1} \X^T \Y & = \hat \Beta  \\
\left[\X^T \X \right]^{-1} \X^T [\X \Beta + {\bf\epsilon}] & = \hat \Beta \\
\left[\X^T \X \right]^{-1} \X^T \X \Beta + \left[\X^T \X \right]^{-1} \X^T {\bf\epsilon} & = \hat \Beta \\
\Beta + \left[\X^T \X \right]^{-1} \X^T {\bf\epsilon} & = \hat \Beta
\end{align*}

}
%
%\frame{
%
%\begin{align*}
%\left[\X^T \X \right]^{-1} \X^T \Y & = \hat \beta  \\
%\left[\X^T \X \right]^{-1} \X^T [\X \beta + {\bf\epsilon}] & = \hat \beta \\
%\left[\X^T \X \right]^{-1} \X^T \X \beta + \left[\X^T \X \right]^{-1} \X^T {\bf\epsilon} & = \hat \beta \\
%\beta + \left[\X^T \X \right]^{-1} \X^T {\bf\epsilon} & = \hat \beta
%\end{align*}
%
%}

\frame{

Thus:

\begin{align*}
 E\left[ V\left[\hat \Beta |\X\right] \right] &=  E\left[ V\left[\Beta + \left[\X^T \X \right]^{-1} \X^T {\bf\epsilon}|\X\right]\right]  \\
&=  E\left[V\left[\left[\X^T \X \right]^{-1} \X^T {\bf\epsilon}|\X\right] \right] \\
& =  E\left[\left[\X^T \X \right]^{-1}  \X^T
V[\epsilon|\X] 
  \X \left[\X^T \X \right]^{-1}\right] 
\end{align*}


}


%\frame{
%
%By Slutsky's Theorem,
%
%\begin{align*}
%n V\left[\hat \Beta\right] &\arrowp n V\left[\left[\E[\X^T \X] \right]^{-1} \X^T {\bf\epsilon}\right] \\
%& = n \left[\E[\X^T \X] \right]^{-1} V\left[ \X^T {\bf\epsilon}\right] \left[\E[\X^T \X] \right]^{-1} \\
%& = n \left[\E[\X^T \X] \right]^{-1} \E\left[ \X^T \text{diag}[{\bf\epsilon}^2] \X \right] \left[\E[\X^T \X] \right]^{-1}
%\end{align*}
%
%The idea is that we can characterize the sampling variability of the estimator just by characterizing the relationship between the errors and the predictors.
%
%}

\frame{

Estimation: plug in residuals. This is the ``sandwich estimator."

\begin{align*}
\hat V\left[\hat \Beta\right] & = \left[\X^T \X \right]^{-1}  \X^T
 {\text{diag}} \left[ (\Y - \X\hat\Beta)^2 \right]
  \X \left[\X^T \X \right]^{-1}.
\end{align*}

This uses the residuals to {\it estimate} the relationship between the $X$ variables and the errors.

\vspace{0.2in}

%N.b., the errors are {\it in expectation} zero. Analogously, the residuals are by definition {\it on average} zero. Otherwise we could improve the fit by changing the constant. C.f., the MMSE approximation.

}

\frame{

One note: there are many variants of the variance estimators, that involve premultiplying by, e.g., $\frac{n}{n-K}$ or $\frac{n}{n-K-2}$, etc. These are analogous to the $\frac{1}{n-1}$ in the sample variance. It's not so simple with the robust variance estimators (many variants: HC0, HC1, HC2) but all are consistent for the SE.
}


\frame{

Many software packages use an alternative. ``Classical" or ``Homoskedastic'' variance estimation. If we {\it assume} that $V[\epsilon | X_1, X_2, ..., X_K] = \sigma^2$ for all values of $X_1, X_2, ..., X_K$: homoskedasticity. Then:

\begin{align*}
V\left[\hat \Beta\right] &= \left[\E[\X^T \X] \right]^{-1} \E\left[ \X^T \text{diag}[{\bf\epsilon}^2] \X \right] \left[\E[\X^T \X] \right]^{-1} \\
&=  \sigma^2 \left[\E[\X^T \X] \right]^{-1} \E\left[ \X^T \I \X \right] \left[\E[\X^T \X] \right]^{-1} \\
&=   \sigma^2 \left[\E[\X^T \X] \right]^{-1} \E\left[ \X^T \X \right] \left[\E[\X^T \X] \right]^{-1} \\
&=   \sigma^2 \left[\E[\X^T \X] \right]^{-1}
 \end{align*}

Where $ \sigma^2 = V [\epsilon] = \E[\epsilon^2]$.

}

%\frame{
%
%There is an alternative. ``Classical" variance estimation. If we {\it assume} that $V[\epsilon | X_1, X_2, ..., X_K] = \sigma^2$ for all values of $X_1, X_2, ..., X_K$: homoskedasticity. Then:
%
%\begin{align*}
%n V\left[\hat \beta\right] &\arrowp  n \left[\E[\X^T \X] \right]^{-1} \E\left[ \X^T {\bf\epsilon} {\bf\epsilon}^T \X \right] \left[\E[\X^T \X] \right]^{-1} \\
%&= n  \sigma^2 \left[\E[\X^T \X] \right]^{-1} \E\left[ \X^T \X \right] \left[\E[\X^T \X] \right]^{-1} \\
%&= n  \sigma^2 \left[\E[\X^T \X] \right]^{-1}
% \end{align*}
%
%Where $ \sigma^2 = V [\epsilon] = \E[\epsilon^2]$.
%
%}

%\frame{
%
%There is an alternative. ``Classical" variance estimation. If we {\it assume} that $V[\epsilon | X_1, X_2, ..., X_K] = \sigma^2$ for all values of $X_1, X_2, ..., X_K$: homoskedasticity. Then:
%
%\begin{align*}
%\hat V\left[\hat \Beta\right] &= \hat \sigma^2 \left[\X^T \X \right]^{-1}.
% \end{align*}
%
%Where $\hat \sigma^2 = \hat V [\Y - \X\hat\Beta]$.
%
%}


\frame{

Classical standard errors make an additional assumption: the magnitude of the errors is {\it unrelated} to the values of the $X$s.

\vspace{0.1in}

This assumption is unlikely to hold. Usually -- in real life -- the errors will be ``heteroskedastic," because in real life, data is messy. If homoskedasticity {\it does} hold, then in large samples, the classical and ``robust" variance estimates will agree.

\vspace{0.1in}

With large $n$ and i.i.d. sampling, always use robust standard errors or bootstrapping to characterize the uncertainty of your regression fit.

%\vspace{0.1in}
%
%What do constant treatment effects imply about homo/heteroskedasticity? What kind of test does the contrapositive allow for?

}


\frame{
If $Y$ and $X$ have been centered and $x$ is a nonstochastic scalar (everything works basically the same if we relax non-stochastic assumption)

\begin{align*}
 V\left[\hat \Beta\right] &= V\left[\Beta +\left[    \frac{\sum_i x_i {\bf\epsilon}_i}{\sum_i x_i^2}\right] \right] \\
 &=    \frac{ V\left[\sum_i x_i {\bf\epsilon}_i\right] }{\left[\sum_i x_i^2\right]^2 } \\
%&=E_x V\left[\left[ \frac{ X_i {\bf\epsilon}_i}{\sum_i X_i^2}\right]
\end{align*}



How does this further simplify if homoskedasticity?

%\vsf
%
%What are the plug-in estimators for these?

}

%\frame{
%\frametitle{Adding variables}
%
%On board draw sampling distribution tables and formulas:
%
%\begin{itemize}
%\item one explanatory var
%\item two explanatory vars
%\end{itemize}
%}








\frame{
\frametitle{Standard Errors/Variance (with homoskedasticity)}

With homoskedasticity, the simple regression model can be written as:

\begin{align*}
 y_i &= \beta_0^s + \beta_1^s x_i + \epsilon_i^s \\
 V&[\epsilon_i]  = \sigma^2_s
\end{align*}

and the two covariate multiple regression model can be written as:

\begin{align*}
 y_i &= \beta_0^m + \beta_1^m x_i + \beta_1^m z_i + \epsilon_i^m \\
 V&[\epsilon_i^m]  = \sigma^2_m
\end{align*}


}


\frame{
\frametitle{Comparing variances (with homoskedasticity)}


\begin{align*}
V[\hat \beta_1^s] &= \frac{\sigma^2_s}{\sum_{i=1}^n (x_i - \bar x)^2 }  \\
V[\hat \beta_1^m] &= \frac{1}{1-R^2_{x \sim z}} \cdot \frac{\sigma^2_m}{\sum_{i=1}^n (x_i - \bar x)^2 } 
\end{align*} \pause

\bigskip

If $x$ is randomized, then what does that tell us about

\bi
\item $\beta_1^s$ and $\beta_1^m$
\item $V[\hat \beta_1^s]$ and $V[\hat \beta_1^m]$
\ei


}



\end{document}